% Generated by scripts/assemble.py; corrections belong in source/corrections.json.
\examhead{45}{上海大学 2024 年数学分析试卷}{47}
% p048-b08; source PDF page 48
\par\Needspace{6\baselineskip}
\blockanchor{p048-b08}
\question{1}(10 分) 叙述函数极限 $\lim_{x \to x_0^-} f(x) = A \in \mathbb{R}$ 的“$\varepsilon - \delta$”定义. 另外, 利用函数极限的定义来证明\par
\nopagebreak[4]
% p048-b09; source PDF page 48
\blockanchor{p048-b09}
\displayformula{\lim _ {x \to 2 ^ {-}} \frac {x ^ {2} - 4}{x} = 0.}
% p048-b10; source PDF page 48
\blockanchor{p048-b10}
\question{2}(10 分) 假设数列 $\{a_{n}\}$ 满足 $\lim_{n\to \infty}\left|\frac{a_n}{a_{n + 1}}\right| = \lambda > 1$. 证明: $\lim_{n\to \infty}a_n = 0$.\par
% p048-b11; source PDF page 48
\blockanchor{p048-b11}
\question{3}(10 分) 计算极限 $\lim_{n\to \infty}\left(\cos \frac{1}{\sqrt{n}}\right)^n$ 与 $\lim_{x\to 0}\frac{\sqrt{1 + \tan x} - \sqrt{1 + \sin x}}{x\ln(1 + x) - x^2}.$\par
% p048-b12; source PDF page 48
\blockanchor{p048-b12}
\question{4}(10 分) 讨论函数 $y = \sin (\sin x)$ 与 $y = x\cos x$ 在 $[0, +\infty)$ 上的一致连续性并给出相应证明.\par
% p048-b13; source PDF page 48
\blockanchor{p048-b13}
\question{5}(10 分) 设函数 $f(x) \in C[0,1]$, 且 $4\int_{0}^{1} f(u) \mathrm{d}u = 1$, 证明: 存在 $\theta \in [0,1]$, 使得 $f(\theta) = \theta^3$.\par
% p048-b14; source PDF page 48
\blockanchor{p048-b14}
\question{6}(10 分) 求幂级数 $\sum_{n=0}^{\infty} (-1)^n \frac{x^n}{n+1}$ 的收敛半径, 收敛域与和函数, 并给出 $\sum_{n=0}^{\infty} (-1)^n \frac{1}{(n+1)2^n}$ 的和.\par
% p048-b15; source PDF page 48
\blockanchor{p048-b15}
\question{7}(10 分) 求空间曲线 $\left\{ \begin{array}{l} x^{2} + y^{2} = 2z; \\ x^{2} + y^{2} + xy = 1. \end{array} \right.$ 上的点到 $xOy$ 平面的最大与最小距离.\par
% p048-b16; source PDF page 48
\blockanchor{p048-b16}
\question{8}(10 分) 设 $A, B$ 为有界数集且 $A \cap B \neq \emptyset$. 证明: $\sup(A \cap B) \leq \min\{\sup A, \sup B\}$, 并给出等号不成立的例子.\par
% p048-b17; source PDF page 48
\blockanchor{p048-b17}
\question{9}(10 分) 设 $f(x) \geq 0$ 在 $[0, +\infty)$ 上可微, $f(0) = 0$, $|f'(x)| \leq f(x)$. 证明: 在 $[0, +\infty)$ 上 $f(x) \equiv 0$.\par
% p048-b18; source PDF page 48
\blockanchor{p048-b18}
\question{10}(10 分) 讨论函数 $z = f(x, y) = \begin{cases} \frac{xy^3}{x^2 + y^4}, & x^2 + y^2 \neq 0; \\ 0, & x^2 + y^2 = 0. \end{cases}$ 在原点 $(0, 0)$ 的连续性, 可偏导性及可微性, 并给出证明.\par
% p048-b19; source PDF page 48
\blockanchor{p048-b19}
\question{11}(10 分) 讨论反常积分 $\int_{2}^{+\infty} \frac{1}{x \ln^2 x} \mathrm{d}x$ 与 $\int_{1}^{+\infty} \frac{1}{x^2 \sqrt{\ln x}} \mathrm{d}x$ 的敛散性并给出相应证明.\par
% p048-b20; source PDF page 48
\blockanchor{p048-b20}
\question{12}(10 分) 将函数 $f(x) = 1 - x^2$ ($0 \leq x \leq \pi$) 展开成余弦级数, 并求级数 $\sum_{n=1}^{\infty} \frac{(-1)^{n-1}}{n^2}$ 的和.\par
% p048-b21; source PDF page 48
\par\Needspace{6\baselineskip}
\blockanchor{p048-b21}
\question{13}(10 分) 求第二类曲面积分\par
\nopagebreak[4]
% p048-b22; source PDF page 48
\blockanchor{p048-b22}
\displayformula{\iint_ {S} x ^ {2} \mathrm{d} y \mathrm{d} z + y ^ {2} \mathrm{d} z \mathrm{d} x + z ^ {2} \mathrm{d} x \mathrm{d} y.}
% p048-b23; source PDF page 48
\blockanchor{p048-b23}
其中 $S$ 是球面 $(x - a)^2 + (y - b)^2 + (z - c)^2 = R^2$，并取外侧.\par
% p049-b00; source PDF page 49
\par\Needspace{6\baselineskip}
\blockanchor{p049-b00}
\question{14}(10 分) 计算曲线积分\par
\nopagebreak[4]
% p049-b01; source PDF page 49
\blockanchor{p049-b01}
\displayformula{I = \int_ {L} (e ^ {x} \sin y - 2 x - 2 y) \mathrm{d} x + (e ^ {x} \cos y - x) \mathrm{d} y.}
% p049-b02; source PDF page 49
\blockanchor{p049-b02}
其中 $L$ 为从 $O(0,0)$ 沿曲线 $y = \sqrt{2x - x^2}$ 到 $A(2,0)$ 的一段有向弧.\par
% p049-b03; source PDF page 49
\blockanchor{p049-b03}
\question{15}(10 分) 假设 $f(x)$ 在 $(- \infty, -1]$ 上连续且 $\int_{-\infty}^{-1} f(x) \mathrm{d}x$ 收敛, 证明: 存在数列 $\{x_n\} \subset (-\infty, -1]$, 使得满足 $\lim_{n \to \infty} x_n = -\infty$ 且 $\lim_{n \to \infty} f(x_n) = 0$. 另外, 假设函数 $g(x): (0, +\infty) \to (0, +\infty)$ 可导, 证明: 存在趋于正无穷大的正数列 $\{x_n\}$, 使得 $g'(x_n) < g(2x_n), n = 1, 2, \cdots$.\par
